\section{Results}
\subsection{The JR column generates and resonantly amplifies the beat}
An input containing only the two tones $f_c\mp\Delta f/2$ produces an output
oscillation at their $11.1$~Hz beat, near the column's alpha resonance
(Fig.~\ref{fig:bif}). The alpha peak is absent from the input spectrum
(Fig.~\ref{fig:demod}).

\begin{figure}[t!]
\centering
\includegraphics[width=\textwidth]{fig_demodulation_twotone.pdf}
\caption{\textbf{The JR column generates a beat absent from the input.}
\emph{(a)} Two equal tones at $94.45$ and $105.55$~Hz and their envelope,
$2\varepsilon|\cos(\Omega t/2)|$. \emph{(b)} Input spectrum: only the two
tones are present. The alpha-band floor is $1.6\times10^{-8}$ of their
amplitude, from finite-window leakage. \emph{(c,d)} The synaptic voltage
$v=y_1-y_2$ and its spectrum show the generated $11.1$~Hz beat.
The field-free column is a stable focus near its alpha Hopf.}
\label{fig:demod}
\end{figure}

In the stable-focus regime ($p>p_H$), the response is approximately Lorentzian
and its peak grows from $0.23$ to $0.43$ to $0.70$~mV toward the Hopf
(Fig.~\ref{fig:res}), as expected from the
$1/\gamma$ susceptibility. Without a drive the focus is silent at $\Omega$;
the nonlinearity generates that component and recurrence amplifies it.
The peak shifts with the driven operating state; the near-Hopf law gives
its leading gain, not its exact frequency. The coupling sweep tests that law
over a decade of damping.
On the alpha branch below $p_H$, the column free-runs at $f_0$, so the lock-in
can contain both driven and autonomous activity. The dominant output frequency
$f_{\rm out}$ then identifies a frequency-dominance region
(Fig.~\ref{fig:entrain}), without establishing phase capture.

\begin{figure}[t!]
\centering
\includegraphics[width=\textwidth]{fig_resonance.pdf}
\caption{\textbf{The JR beat resonance sharpens toward the Hopf.}
Stable-focus response at $\varepsilon=0.3$~mV. \emph{(a)} Lock-in amplitude
against beat frequency; the finite-amplitude peak moves from $11.08$ to
$10.93$~Hz across the curves. \emph{(b)} Peak amplitude increases as $p$
approaches the upper Hopf from above. These three points show sharpening
and growth; the coupling sweep quantifies the $1/\gamma$ scaling and
its saturation.}
\label{fig:res}
\end{figure}

\subsection{Carrier independence and the kHz roll-off}
At fixed post-access polarization, the open-loop detector and near-Hopf
JR column are carrier-independent outside their natural band
(Fig.~\ref{fig:carrier}). At fixed applied field, the point-membrane
response instead falls as $1/f_c^2$ above its corner
(Fig.~\ref{fig:khz}); the sealed-cable estimate is compared with the
point-membrane cases in Fig.~\ref{fig:feasibility}a.

\begin{figure}[t!]
\centering
\includegraphics[width=\textwidth]{fig_carrier_independence.pdf}
\caption{\textbf{JR carrier independence at fixed polarization.}
\emph{(a)} Open- and closed-loop beat responses over $f_c=25$--$300$~Hz,
with deviations near the natural band. Each is normalized by its own
$300$~Hz response: $0.00670$~s$^{-1}$ open loop and $0.512$~mV closed
loop, giving a rate-to-voltage transfer of $76.4$~mV\,s.
\emph{(b)} Relative synaptic transfer $|H_s(\omega)|/|H_s(0)|$: about
$0.7$ at $10$~Hz and $2.5\%$ at $100$~Hz. Dotted: synaptic corner
$a/(2\pi)\approx16$~Hz; shaded: $8$--$13$~Hz beat band; dashed: carrier.}
\label{fig:carrier}
\end{figure}

With an explicit first-order access state per modeled compartment before
the sigmoid, the full JR column demodulates and resonates in the assumed
$\tau=0.2$~ms case (Fig.~\ref{fig:khz_direct}a). Genuine tones centered
at $2$~kHz generate the absent $10.914$~Hz beat, with analytic access
factor $|H_m|^2=0.137$ (Fig.~\ref{fig:khz_twotone}). These are consequences
of the assumed access time and coupling; neither is calibrated by the simulation.

Figure~\ref{fig:feasibility} converts this access assumption into a source
budget. With $\varepsilon_j=\lambda_0 E_j|H_m(\omega_j)|$, the equivalent
direct beat input is $\varepsilon_{\rm eq}=Y_\Omega/\Sig'$, before recurrent
gain. At $2$~kHz, $\tau_m=0.2$~ms, $\lambda_0=0.5$~mV per V/m, and
$0.5$~V/m per tone, it has magnitude $1.36~\mu$V; the steady equivalent
is nearly identical. Reducing the assumed coupling to $0.15$ scales these
sources by $0.09$. Thus the matched-polarization $2$--$5\%$ comparison does
not carry over unchanged to matched applied fields.

At equal fields, the $0.2$~ms point-membrane model predicts a narrowband
$9/1$~kHz source ratio of $0.020$. The limit $1/81$ applies only when
both bands lie above their access corners. A sealed long cable instead
approaches a $1/9$ ratio. These calculations locate conditional input budgets;
they do not calibrate the biological transfer.

\begin{figure}[t!]
\centering
\includegraphics[width=\textwidth]{fig_feasibility_khz.pdf}
\caption{\textbf{Carrier access sets the source budget at fixed field.}
\emph{(a)} Exact two-tone access factor for tones $10$~Hz apart; dotted,
sealed-cable approximation at $0.2$~ms.
\emph{(b)} Beat-equivalent direct polarization magnitude at the pinned JR state,
$E=0.3$~V/m per tone and $\lambda_0=0.5$~mV per V/m. For $\lambda_0=0.15$,
multiply by $0.09$.
\emph{(c)} Field dependence at $0.2$~ms for both couplings; star marks
the illustrative budget. All sources are leading-order, before recurrent
gain; $\lambda_0$ and the access time are assumptions.}
\label{fig:feasibility}
\end{figure}

\subsection{Sigmoid curvature predicts the JR beat source}\label{sec:results-curvature}
The signed open-loop beat follows $\half\Sig''(v^*)\varepsilon^2m$, vanishing
and reversing phase at the inflection (Fig.~\ref{fig:opp}). The closed-loop
examples occupy the negative-curvature flank. At matched network
susceptibility, either flank follows the same gain law, with opposite source
phase when the curvature changes sign.

\begin{figure}[t!]
\centering
\includegraphics[width=\textwidth]{fig_curvature_twotone.pdf}
\caption{\textbf{Curvature determines the JR detector.}
\emph{(a)} The signed open-loop beat follows
$\Sig''(v^*)\varepsilon^2/2$ within $0.9\%$, vanishing and reversing sign
at $v_0$. \emph{(b)} The recurrent response has amplitude slope $1.99$.
\emph{(c)} Linearizing the receiving sigmoid reduces the beat $634$-fold.
\emph{(d)} Open-loop two-tone/AM beat ratio at $v^*=8$~mV:
$1.0000$ at $0.02$~mV and $0.988$ at $1$~mV. This checks the leading
local source, not the nonlinear recurrent response near the Hopf.}
\label{fig:opp}
\end{figure}

In the recurrent column, the response scales quadratically at small drive
(log--log slope $1.99$, measured as in \S\ref{sec:transduction-methods}),
and linearizing the field-receiving sigmoid reduces it
from $1.34$ to $0.0021$~mV (Fig.~\ref{fig:opp}b,c). The open-loop surrogate
comparison in panel~(d) supports its use for leading-order beat sweeps, with
the waveform and near-Hopf limitations specified in \S\ref{sec:waveforms}.

\subsection{Realistic-head comparison of envelope and quadratic maps}\label{sec:fem-results}
In this single-subject geometric illustration, changing the current ratio
moves the envelope and mean-square ROI centroids in opposite directions,
while the quadratic beat map only rescales. Each centroid uses the top 10\%
of hippocampal nodes; it is not a whole-brain maximum.
Figure~\ref{fig:fem} compares the maps of Table~\ref{tab:transduction-maps}
for the Violante montage in IXI002, using the same polarizations and fixed
1:1 projection (\S\ref{sec:fem-methods}).

\begin{table}[htbp]
\centering\small
\caption{\textbf{Hippocampal maps under current-ratio changes.}
Contrasts divide the ROI p95 or raw mean by the reference p95; parentheses
give square-root, degree-matched p95 contrasts. Focus shifts use the top
10\% of ROI nodes, anterior positive. Indented rows restrict the contrasts
to gray matter (GM); other entries use the gray- and white-matter ROI.}
\label{tab:fem-summary}
\setlength{\tabcolsep}{4pt}
\begin{tabular}{@{}lrrrrrr@{}}
\toprule
Map & \shortstack{ROI mean\\1:1} & \shortstack{p95 contrast\\1:1} &
\shortstack{p95 contrast\\1:3} & \shortstack{Mean contrast\\1:1 / 1:3} &
\shortstack{Mean\\change} & \shortstack{Focus\\shift}\\
\midrule
Envelope & $0.16$ mV & $1.22$ & $0.95$ & $0.81\,/\,0.58$ & $-45\%$ & $+2.5$ mm\\
\quad GM only & & $1.09$ & $0.84$ & $0.76\,/\,0.55$ & & \\
Quadratic beat & $0.018$ mV$^2$ & $1.17$ ($1.08$) & $1.17$ ($1.08$) & $0.61\,/\,0.61$ & $-25\%$ & $0$ mm\\
\quad GM only & & $0.94$ ($0.97$) & $0.94$ ($0.97$) & $0.51\,/\,0.51$ & & \\
Mean square & $0.0094$ mV$^2$ & $0.92$ ($0.96$) & $0.90$ ($0.95$) & $0.47\,/\,0.54$ & $+24\%$ & $-4.5$ mm\\
\quad GM only & & $0.71$ ($0.85$) & $0.76$ ($0.87$) & $0.40\,/\,0.46$ & & \\
\bottomrule
\end{tabular}
\end{table}

\begin{figure}[t!]
\centering
\includegraphics[width=\textwidth]{fig_fem.pdf}
\caption{\textbf{Envelope and quadratic maps in IXI002.}
Violante montage; left hippocampus outlined in gold, anterior to the left.
\emph{(a--d)} Sagittal maps at $2+2$~mA (top) and $1+3$~mA (bottom):
envelope excursion, quadratic excursion, mean-square polarization, and
isotropic beat. The first three use the fixed 1:1 envelope-maximizing
direction. Each color scale ends at its 1:1 gray-matter 99th percentile.
Circles and triangles mark the projected centroids of the top 10\% of ROI
nodes at 1:1 and 1:3, not whole-brain maxima.
\emph{(e)} Target-to-reference contrast; dashed curves use square-root
maps to match degree. \emph{(f)} Anterior focus displacement: top 10\%
(line), range over 5--50\% (band), colors as in (e). \emph{(g)} Contrasts at 1:1, raw
(solid) and degree-matched (hatched).
The reference is gray matter outside the ROI. All maps precede detector
gain: $\mathcal I$ weights $Y_{\rm DC}=\Sig''\mathcal I/2$, but does not set its sign.}
\label{fig:fem}
\end{figure}

The pair-current product falls from $4$ to $3$~mA$^2$, so both beat maps fall to $75\%$ at every node, to machine precision, with unchanged normalized patterns, contrasts, and foci. The envelope focus moves toward the pair whose current falls; the mean-square focus moves toward the pair whose current rises. For the contrast comparison at balance, degree matching reduces the quadratic contrast to $1.08$, below the envelope's $1.22$; the isotropic beat gives $1.18$ ($1.08$ degree-matched). The envelope's anterolateral tail above the reference p95 is 52\% white matter, against 25\% in the ROI. Restricting the ROI to gray matter removes the beat preference: only the envelope exceeds one at 1:1, and no map does at 1:3 (Table~\ref{tab:fem-summary}). No mean-based contrast exceeds one. These geometric contrasts do not establish physiological deep selectivity, which also depends on detector and circuit weighting.

The envelope contrast peaks at $1.24$ near balance. Once one carrier becomes
weaker across most of the ROI, further imbalance changes amplitude more than
focus; the mean-square focus moves continuously. Across all 32 combinations
of montage, label, tissue definition, and focus fraction, the envelope shift
remains anterior, the mean-square shift posterior, and the quadratic shift
zero. Radial projections preserve these directions but change contrasts more
than the atlas label does (supplement, \S\ref{supp:fem-sensitivities}).

A laminar-aligned detector samples median fractions $0.29$, $0.11$, and $0.12$ of the envelope, quadratic, and mean-square maps' respective orientation maxima. Evaluated on 934 of 1,568 ROI nodes, these fractions quantify the anatomical access hypothesis (supplement, \S\ref{supp:fem-orientation}). The median ratios of projected maps across the tissue interfaces are $1.13$--$1.34$; mesh convergence remains untested (supplement, \S\ref{supp:fem-interfaces}).
\ifjne\else
\subsection{Sensitivity of focus displacement}\label{supp:fem-sensitivities}
The spatial ordering in the article persists across montage, labeling, tissue,
and focus-fraction choices. Table~\ref{tab:fem-sensitivities} gives the ranges;
orientation changes absolute contrast more than the atlas label in this subject.
For the reference projection and label, the top-5--50\% focus moves
$2.2$--$4.0$~mm anteriorly for the envelope and $2.4$--$8.2$~mm posteriorly
for mean-square polarization.

The sweep shows two kinds of steering (Fig.~\ref{fig:fem}e,f). The envelope ratio peaks
near balance ($1.24$ at a posterior share of $0.525$) and falls to $0.89$--$0.94$ at
the extremes. Its focus moves only while the two carriers are comparable inside the
hippocampus. At a posterior share of $0.6$ or more, the anterior pair is the weaker
carrier at more than 96\% of hippocampal nodes; at $0.4$ or less, the posterior pair is. The
envelope map then equals twice the weaker pair's polarization over most of the ROI,
and its measured focus changes little with further imbalance. The focus of $\mathcal I$ moves
continuously instead, from $1.5$~mm anterior at a posterior share of $0.1$ to $6.3$~mm
posterior at $0.9$, while its ratio stays between $0.79$ and $0.94$ ($0.89$--$0.97$
degree-matched). The quadratic excursion keeps one ratio and one focus across the
whole sweep.

\begin{table}[htbp]\centering\small
\caption{Focus displacement from 1:1 to 1:3. Envelope shifts are anterior;
mean-square shifts are posterior. The quadratic beat shift is zero in every case.
Contrasts quoted in the last column are raw 1:1 envelope/quadratic contrasts.}
\label{tab:fem-sensitivities}
\begin{tabular}{@{}lrrl@{}}
\toprule
Sensitivity & Envelope (mm) & Mean square (mm) & Contrasts\\
\midrule
Atlas label & $2.0$--$3.8$ & $3.2$--$5.9$ & $1.17$ / $1.04$\\
Radial direction & $1.9$--$2.8$ & $2.0$--$3.7$ & $1.62$ / $1.97$\\
All 32 combinations & $0.7$--$8.1$ & $0.5$--$10.6$ & \\
\bottomrule
\end{tabular}
\end{table}


The gray-matter-only comparison uses 1,174 FreeSurfer hippocampal nodes
against the unchanged gray-matter reference. Table~\ref{tab:fem-gm-contrasts}
gives both contrast definitions for all four maps. Degree matching takes
the square root of a map before computing its statistics. At 1:1 only the
envelope p95 contrast exceeds one; at 1:3 none does. No raw or degree-matched
mean contrast reaches one in either tissue definition.

\begin{table}[htbp]\centering\small
\caption{Gray-matter-only hippocampal contrasts, 1:1 / 1:3. Each denominator
is the reference p95 of the same raw or square-root map.}
\label{tab:fem-gm-contrasts}
\setlength{\tabcolsep}{4pt}
\begin{tabular}{@{}lrrrr@{}}
\toprule
& \multicolumn{2}{c}{Raw map} & \multicolumn{2}{c}{Square-root map}\\
Map & p95 contrast & Mean contrast & p95 contrast & Mean contrast\\
\midrule
Envelope & $1.086/0.839$ & $0.757/0.547$ & -- & --\\
Quadratic beat & $0.935/0.935$ & $0.508/0.508$ & $0.967/0.967$ & $0.696/0.696$\\
Isotropic beat & $0.943/0.943$ & $0.512/0.512$ & $0.971/0.971$ & $0.698/0.698$\\
Mean square & $0.715/0.765$ & $0.396/0.462$ & $0.845/0.874$ & $0.615/0.665$\\
\bottomrule
\end{tabular}
\end{table}

\subsection{Field orientation relative to hippocampal anatomy}\label{supp:fem-orientation}
The fields in IXI002 follow neither the hippocampal long axis nor the local
laminar direction. Table~\ref{tab:fem-orientation} compares the carrier fields
$\mathbf E_1,\mathbf E_2$ and each montage's balanced-envelope-maximizing
direction $\mathbf n^*$ at $2+2$~mA.
The principal coordinate axes of the 1,568-node FreeSurfer hippocampal region
of interest (ROI) are PC1, the anterior-positive long axis; PC2, lateral-positive; and
PC3$={\rm PC1}\times{\rm PC2}$, normal to its best-fit plane. Angles are
unsigned, from $0$ to $90^\circ$, and nodes receive equal weight.
Weighting by field magnitude, or envelope maximum for $\mathbf n^*$,
changes these angular medians by at most about $1.5^\circ$.

\begin{table}[htbp]
\centering\small
\setlength{\tabcolsep}{3pt}
\begin{tabular}{@{}llrrrr@{}}
\toprule
Montage & Direction & PC1 & PC2 & PC3 & Laminar\\
\midrule
Violante & $\mathbf E_1$ & 71 (63--78) & 25 (17--33) & 79 (73--84) & 69 (53--80)\\
 & $\mathbf E_2$ & 58 (51--67) & 35 (27--42) & 83 (77--87) & 72 (57--81)\\
 & $\mathbf n^*$ & 66 (62--71) & 28 (22--34) & 80 (73--85) & 69 (53--80)\\
Yatsuda & $\mathbf E_1$ & 71 (64--78) & 25 (20--35) & 76 (69--84) & 74 (55--83)\\
 & $\mathbf E_2$ & 77 (71--83) & 21 (15--27) & 79 (70--85) & 74 (54--83)\\
 & $\mathbf n^*$ & 75 (66--83) & 23 (16--32) & 78 (70--85) & 74 (56--83)\\
\bottomrule
\end{tabular}
\caption{Angles in degrees, median (interquartile range). Coordinate-axis
comparisons use all 1,568 ROI nodes; laminar comparisons use the 934 nodes
where that direction is defined. Yatsuda denotes the FT7--F8, TP7--P8
rendering.}
\label{tab:fem-orientation}
\end{table}

The laminar direction is the normalized gradient of HippUnfold's inner--outer
Laplace coordinate \cite{dekraker2022}, sampled at the nearest voxel without
smoothing. HippUnfold 1.5.2-pre.2 processed the same T1 image; its container,
arguments and segmentation-model hash are recorded in
\path{docs/fem_ixi002_violante.md}. The coordinate has a $0.2$~mm grid
in T1 world coordinates. It defines directions at 934 nodes (59.6\% of the
ROI); 565 are interior, with their voxel and all six face neighbors inside
the Laplace domain. The laminar direction lies $33^\circ$ (21--54) from PC3
and $76^\circ$ (55--83) from PC2. Its median spread about its own principal
axis is $31^\circ$ (23--51), so one uniform direction cannot follow it
throughout the region.

A laminar-aligned detector samples smaller polarizations than an optimally
oriented one. In the Violante montage, median absolute projections of the
unit carrier directions onto the laminar direction are 0.36 and 0.31.
The map fractions in Table~\ref{tab:fem-laminar-fractions} are computed
node by node, each divided by that map's own maximum over orientation.
The quadratic and mean-square denominators therefore need not occur at
$\mathbf n^*$. These fractions describe geometry under a laminar-detector
hypothesis, before state-dependent detection or recurrent gain.

\begin{table}[htbp]
\centering\small
\begin{tabular}{@{}llrrr@{}}
\toprule
Montage & Nodes & Envelope & Quadratic beat & Mean square\\
\midrule
Violante & All defined (934) & 0.29 (0.11--0.57) & 0.11 (0.02--0.34) & 0.12 (0.04--0.33)\\
 & Interior (565) & 0.25 (0.10--0.54) & 0.08 (0.02--0.29) & 0.09 (0.03--0.31)\\
Yatsuda & All defined (934) & 0.25 (0.09--0.55) & 0.07 (0.01--0.32) & 0.08 (0.02--0.33)\\
 & Interior (565) & 0.25 (0.10--0.47) & 0.07 (0.02--0.25) & 0.08 (0.02--0.26)\\
\bottomrule
\end{tabular}
\caption{Laminar-projected map divided by its orientation maximum, median
(interquartile range). Restricting the Violante comparison to interior nodes
gives carrier and $\mathbf n^*$ angles to the laminar direction of
$71^\circ$, $74^\circ$, and $71^\circ$, respectively.}
\label{tab:fem-laminar-fractions}
\end{table}

The carrier directions themselves differ by a median $14.2^\circ$ (10.2--18.6) in the Violante montage and $7.5^\circ$ (4.2--11.4) in the Yatsuda rendering; their 95th percentiles are $29.2^\circ$ and $16.8^\circ$. The renderings also differ in electrode layout. With the tested 20-mm discs, the Violante e2 center is 9.4 and 18.4~mm from AF8 and Fp2, and e4 is 10.4 and 19.0~mm from P8 and PO8. These separations meet the 20-mm overlap criterion if those cap positions are occupied, for this layout and disc size.

\subsection{Continuum bounds and recovered interface jumps}\label{supp:fem-interfaces}
Median map jumps across tissue interfaces are 1.13--1.34 along the common
direction $\mathbf n^*$ (Table~\ref{tab:fem-interface-jumps}). These
projected ratios differ from the conductivity bounds on field magnitude
and maps maximized over orientation. All measurements here use IXI002,
the Violante montage, and $2+2$~mA.

Interfaces are faces shared by tetrahedra of different tissue tags:
241,780 CSF/gray-matter faces and 253,608 gray/white-matter faces, with no
white-matter/CSF faces. The median CSF/gray-matter edge is 1.51~mm
(interquartile range 1.37--1.66). At each interface node, each side's field
is averaged over its incident interface-face element samples; the normal
is the normalized sum of their area vectors. Both sides use the
balanced-envelope-maximizing direction $\mathbf n^*$ of the recovered
gray-matter nodal fields. This direction is not reoptimized for the
side-averaged fields. Fixed-projection ratios retain high-side values
exceeding $10^{-12}$ times that side's map maximum.

For a continuum interface between isotropic conductivities, let
$r=\sigma_{\rm high}/\sigma_{\rm low}\geq1$, where ``high'' and ``low''
identify the more and less conductive sides, and let $\boldsymbol\nu$ be
the unit interface normal. Tangential-field and normal-current continuity give
\begin{equation}
\mathbf E_{j,\rm low}=A\mathbf E_{j,\rm high},\qquad
A=\mathbf 1+(r-1)\boldsymbol\nu\boldsymbol\nu^{\mathsf T},\qquad j=1,2,
\label{eq:fem-interface-transform}
\end{equation}
where $\mathbf 1$ is the identity. The symmetric matrix $A$ has eigenvalues
$1,1,r$. For unit $\mathbf n$, define
$\mathbf m=A\mathbf n/|A\mathbf n|$; then
$\mathbf n\cdot A\mathbf E_j=|A\mathbf n|\,\mathbf m\cdot\mathbf E_j$,
with $1\leq|A\mathbf n|\leq r$.
For the map bounds, the scalar field-to-polarization coupling is the same
on both sides. Let $M_k$ denote a map's maximum over unit orientation, with degree $k=1$
for envelope excursion and $k=2$ for quadratic excursion or mean-square
polarization. For nonzero $M_{k,\rm high}$,
\begin{equation}
1\leq\frac{M_{k,\rm low}}{M_{k,\rm high}}\leq r^k.
\label{eq:fem-interface-bound}
\end{equation}
The upper bound follows by maximizing the projection identity. For the lower
bound, choose $\mathbf n$ parallel to $A^{-1}\mathbf m$ for a high-side
maximizer $\mathbf m$; its scale factor is at least one. Field magnitude
obeys the corresponding degree-one bound. Here $r=4.475$ at CSF/gray matter
and $2.667$ at gray/white matter, giving squared bounds 20.03 and 7.11.

A generic fixed oblique projection has neither the floor one nor the ceiling
$r^k$: normal and tangential contributions can cancel on either side.
An exact low-side envelope maximizer nevertheless retains the envelope floor,
and an exact high-side maximizer retains its ceiling. Gray matter is the
low side at CSF/gray matter and the high side at gray/white matter, but the
measured $\mathbf n^*$ maximizes recovered nodal fields, not the side averages
used for jumps. Optimization and field averaging must therefore be
distinguished when applying the continuum bounds.

\begin{table}[htbp]
\centering\small
\begin{tabular}{@{}lrrrr@{}}
\toprule
Low/high ratio & Nodes & Envelope & Quadratic beat & Mean square\\
\midrule
GM/CSF & 120,892 & 1.17 (1.04--1.53) & 1.28 (1.06--2.07) & 1.25 (1.04--1.97)\\
WM/GM & 126,426 & 1.13 (1.03--1.38) & 1.33 (1.08--2.01) & 1.34 (1.08--2.05)\\
\bottomrule
\end{tabular}
\caption{Projected map jumps along the same $\mathbf n^*$ on both sides,
median (interquartile range). GM denotes gray matter and WM white matter.
The conductivity bounds in Eq.~\eqref{eq:fem-interface-bound} apply to
each map's own orientation maximum, not these projections.}
\label{tab:fem-interface-jumps}
\end{table}

The median angle between $\mathbf n^*$ and the interface normal is $64^\circ$ at
CSF/gray matter and $63^\circ$ at gray/white matter. Across maps and interfaces,
4.4--12.8\% of the projected ratios lie below one and 0.2--2.9\% exceed
$r^k$. The largest envelope and quadratic ratios are $7.2\times10^4$ and
$1.29\times10^5$, respectively; large ratios can arise from small denominators
and do not measure large absolute exposure.

Even separately optimized maps depart from the continuum bounds in the
recovered fields (Table~\ref{tab:fem-interface-tails}): 96.7--98.3\% lie
within them, 1.6--3.3\% below one, and fewer than 0.05\% above the ceiling.
First percentiles are 0.988--0.997, but minima reach 0.686.

\begin{table}[htbp]
\centering\small
\begin{tabular}{@{}llrrrrr@{}}
\toprule
Interface & Map & Below 1 & Minimum & Above $r^k$ & Maximum & p1\\
\midrule
CSF/GM & Envelope & 3513 & 0.885 & 33 & 7.90 & 0.995\\
 & Quadratic beat & 3845 & 0.796 & 3 & 20.63 & 0.989\\
 & Mean square & 3936 & 0.686 & 3 & 21.27 & 0.988\\
GM/WM & Envelope & 2275 & 0.920 & 45 & 3.51 & 0.997\\
 & Quadratic beat & 2075 & 0.877 & 33 & 8.67 & 0.995\\
 & Mean square & 2211 & 0.880 & 34 & 8.57 & 0.995\\
\bottomrule
\end{tabular}
\caption{Ratios of each map's own orientation maxima, low/high conductivity.
Tail entries are node counts, with denominators in Table~\ref{tab:fem-interface-jumps};
p1 denotes the first percentile. Counts use a relative tolerance of $10^{-9}$
at the ceiling and an absolute tolerance of $10^{-9}$ at one.}
\label{tab:fem-interface-tails}
\end{table}

Component measurements quantify these departures without establishing their
cause. The high-side normal fraction of $\mathbf E_1$ has nodal medians
20\% and 46\% at the two interfaces. Applying the continuum transform to
each high-side field predicts median projected jumps of 1.26--1.80.
For faces whose high-side normal component is at least 10\% of field magnitude,
the realized $\mathbf E_1$ normal jump is 2.15 (1.45--3.18) at CSF/gray matter
and 1.88 (1.51--2.24) at gray/white matter. Varying that cutoff from 0 to 50\%
gives medians 1.97--2.17 and 1.88--1.90; the corresponding $\mathbf E_2$
ranges are 1.99--2.17 and 1.86--1.88. Element-face tangential continuity holds
to rounding (largest residual $5.2\times10^{-11}$ of the high-side field).
Nodal side averages do not retain this identity: their median tangential
residuals are about 1\% of the field.

At nodes above the own-maximum ceiling, the median realized $\mathbf E_1$
normal jump exceeds $r$ for all six map/interface combinations; the envelope
subsets give 7.08 and 2.80. Tangential residuals here are normalized by the
high-side tangential magnitude; their medians occupy the 60th--75th
percentiles of all-node residuals. Below-floor subsets have median normal
ratios 0.81--1.14. These conditional medians associate the map tails with
normal-component scatter; they do not determine the cause of every violation.
Continuous linear finite elements enforce the normal-current condition in
the weak formulation, not face by face. Whether refinement closes the gap
remains unresolved; refining this mesh would also retain its segmented boundary.

The exposure statistics compare nodes near interfaces with those farther away. ``Near'' means within 2~mm of an interface vertex or face centroid. Gray-matter 99th percentiles near CSF are 0.90, 0.86, and 0.93 times those farther away for envelope, quadratic, and mean-square maps. White-matter ratios near the gray/white interface are 1.25, 1.62, and 1.57. Of each map's top 1\% of brain nodes, 67--73\% are white matter and 81--82\% lie near the gray/white interface, compared with 73\% of all brain nodes; 24--31\% lie near CSF, compared with 47.5\% of all nodes.

A separate check of local maxima holds the projection along the spatially fixed $x$ direction. A strict local maximum exceeds every node sharing a tetrahedron edge within the same tissue; a boundary node belongs to a face adjoining another tissue. Boundary fractions of mean-square maxima are 98.0\% in gray matter and 96.5\% in white matter, against boundary fractions of 90.0\% and 73.4\% among all nodes. The quadratic beat gives similar fractions, 97.9\% and 96.1\%, despite lacking the subharmonic restriction. Mean-square interior maxima number 262 of 13,144 in gray matter and 227 of 6,395 in white matter; the quadratic counts are 270 of 12,659 and 245 of 6,262. With 95.4\% of all brain nodes already within 2~mm of an interface, these mesh-dependent statistics do not sharply test the continuum restriction in \S\ref{sec:dc}.

\fi

\subsection{JR gain, mean shifts, and the tACS comparison}\label{sec:jcurve}
For the AM surrogate, pinning $v^*=8.30$~mV through $p(C)$ fixes the curvature and open-loop
beat $A_{\rm open}\simeq\half|\Sig''(v^*)|\varepsilon^2m$.
As $C\to C^*\approx137$, damping tends to zero, the response grows about
$30$-fold, and resonance sharpens (Fig.~\ref{fig:jcurve}a,b).
This isolates a network effect: recurrent transfer changes at fixed detector curvature.
The rate-to-voltage gain $A_\Omega/A_{\rm open}$ (mV\,s) follows
$1/\gamma$ before saturating (panel~c). Approaching the Hopf through $p$
gives similar leading scaling. Finite-distance gains can differ because
modal frequencies and input--output projections vary along each path.

\begin{figure}[t!]
\centering
\includegraphics[width=\textwidth]{fig_jcurve.pdf}
\caption{\textbf{JR coupling controls amplification at fixed detection.}
AM surrogate, Eq.~\eqref{eq:amsur}, with $p(C)$ pinning the field-free state,
curvature $\Sig''(v^*)=-0.151$~s$^{-1}$mV$^{-2}$, and source $A_{\rm open}$.
\emph{(a,b)} The response grows and its resonance sharpens as $C$ approaches
$C^*$ on the stable-focus side ($\gamma>0$). \emph{(c)} Rate-to-voltage gain
$A_\Omega/A_{\rm open}$ (mV\,s) follows $1/\gamma$ (dashed) before saturating.}
\label{fig:jcurve}
\end{figure}

\paragraph{Beat modulation and mean rate.}\label{sec:timing}
With genuine two tones at $\Delta f=f_0$, the same pinned sweep
increases modulation of $r_P=\Sig(v+s)$ by $13.2$-fold ($26$ to $339$~mHz),
following $1/\gamma$ before saturation. The same rate's mean shift grows
$4.4$-fold ($-1.8$ to $-7.7$~mHz), remaining below $0.2\%$ of its
$3.9$~Hz baseline (Fig.~\ref{fig:timing}). Their magnitude ratio rises
from $14.5$ to $58$ near $\gamma=1$~s$^{-1}$, then falls to $44$ as
the amplified oscillation rectifies. In this weakly driven regime, a small
mean change accompanies larger beat modulation, consistent with primate observations
\cite{vieira2024}, with both compared in rate units. The surrogate
comparison is supplementary (\S\ref{supp:waveform-checks}).

The detector's steady and beat sources meet different transfers:
$|G(\omega_0)|\propto1/(2\gamma\omega_0)$ grows toward the Hopf, while
$|G(0)|\propto1/(\omega_0^2+\gamma^2)$ stays finite. The amplified beat
also rectifies, adding a mean shift quadratic in its amplitude that can
grow as $1/\gamma^2$ before saturation. A bounded zero-frequency transfer
therefore does not bound the full mean shift.

\paragraph{Comparison with tACS.}\label{sec:tacs-results}
TI and tACS share the linear amplifier, so matching their generated in-band
drives at a common state matches their linear outputs; resonance alone gives
TI no efficiency advantage.
For the AM surrogate with amplitude $\varepsilon_{\rm TI}$ and modulation depth
$m$, and a tACS input of amplitude $\varepsilon_{\rm tACS}$, the ratio of voltage
lock-in amplitudes in linear response is
\[
\mathcal R_{\rm lin}
=\frac{|\Sig''(v^*)|\,\varepsilon_{\rm TI}^2m}
{2\Sig'(v^*)\,\varepsilon_{\rm tACS}}.
\]
The common linear transfer cancels, giving a predicted ratio of $0.143$ for the original unequal-input comparison. Genuine two-tone responses depart from this prediction near the Hopf. When the leading sources are matched instead, the TI/tACS output ratios are $0.90$--$0.99$ at the three tested couplings. The larger direct input saturates earlier and shifts the operating state, which field-free pinning does not cancel. Numerical controls for these unequal-input comparisons are in the supplement (\S\ref{supp:tacs-saturation}). Equal-per-carrier, total-peak, and RMS matching are distinguished in \S\ref{sec:direct-comparison}.

\ifjne\else
\subsection{TI/tACS ratios and finite-amplitude saturation}\label{supp:tacs-saturation}
Table~\ref{tab:waveform-checks} gives the AM and genuine two-tone response
ratios for the three coupling controls specified in
\S\ref{app:jr-sweep-settings}. The predicted linear ratio is $0.143$ for
the unequal amplitudes in Figs.~\ref{fig:jcurve} and~\ref{fig:tacs}; the
common transfer cancels at a fixed state. The stronger direct response
saturates earlier near the Hopf: multiplying its input by $6.99$ increases
its output by only $6.97$, $3.95$, and $1.96$ at those three couplings.
Stimulation-induced mean shifts, which field-free pinning does not cancel,
account for the departure from the linear ratio. These are not equal-RMS
comparisons; matching conventions are given in \S\ref{sec:direct-comparison}.

\fi

\begin{figure}[t!]
\centering
\includegraphics[width=\textwidth]{fig_nmm2_resonance_twotone.pdf}
\caption{\textbf{Gamma resonance in the PING mean field.}
Excitatory-rate lock-in under genuine two-tone input, carrier $300$~Hz.
\emph{(a)} Below $\bar\eta_H=1.110$, forced-resonance peaks near
$f_0\approx54$~Hz grow from $0.0014$ to $0.0030$ to $0.0055$~ms$^{-1}$
toward the Hopf. \emph{(b)} Above it, peaks in the autonomous-gamma regime
shift from $56$ to $65$ to $83$~Hz with increasing $\bar\eta$.
The voltage nonlinearity $v^2$ generates the beat absent from the input.
AM/two-tone comparisons are in the supplement (\S\ref{supp:waveform-checks}).}
\label{fig:nmm2_res}
\end{figure}

\begin{figure}[t!]
\centering
\includegraphics[width=\textwidth]{fig_mpr_crossover.pdf}
\caption{\textbf{Static curvature at slow carriers, slope at fast ones.} \emph{(a)}
Second-order DC rate shift of an uncoupled MPR population ($\Delta=1$) under a
current carrier, Eq.~\eqref{eq:crossover}, for operating points below, at and above the
inflection of $F(Q)$. Dotted: the static limit $\tfrac14F''$; dashed: the fast-carrier
limit $F'/[2(\omega\tau_m)^2]$; circles: RK4 integrations at $I_{\rm ac}=0.05$ averaged over whole carrier
periods (relative error below $10^{-3}$). Above the inflection the sign changes with carrier frequency. The horizontal coordinate is $\omega\tau_m$; the carrier markers
use $\tau_m=\tau_E=15$~ms. \emph{(b)} PING of Eq.~\eqref{eq:nmm2-ping}: relative error of the
fast-carrier difference-frequency source against the exact second-order response at four
stable operating points ($\gamma$ the damping of the gamma mode), two tones centered at
$300$, $1000$ and $2000$~Hz.}
\label{fig:crossover}
\end{figure}

\subsection{PING amplifies the dynamic QIF beat source}\label{sec:nmm2-results}
With genuine $300$~Hz tones entering E current, the PING rate $r_E$
shows forced gamma resonance below the Hopf and a beat response superposed
on autonomous gamma above it (Fig.~\ref{fig:nmm2_res}).
The resonance grows through the Hopf and shifts with state, forming a diagonal
ridge in the $(\Delta f,\bar\eta)$ plane of Fig.~\ref{fig:nmm2_map}, unlike
the JR example \cite{clusella2023}.

At $300$~Hz, voltage excursions mix through $v_E^2$ on the fast side of
the detector's crossover (Fig.~\ref{fig:crossover}a). Each current tone
produces a voltage phasor $I_j/(i\omega_j\tau_E)$, giving a beat source
$I_1I_2/(\tau_E^2\omega_1\omega_2)$ in the equation for $\tau_E\dot v_E$.
After division by $\tau_E$ and propagation through the network transfer at
$\Omega$, it matches the full second-order response at four stable operating
points within $1.1\%$ at $300$~Hz, $0.1\%$ at $1$~kHz, and $0.03\%$
at $2$~kHz (Fig.~\ref{fig:crossover}b).

The unpinned PING coupling sweep also approaches a gamma Hopf, but two
modes compete for the response peak. Near the Hopf, the critical mode's
linear gain has a fitted damping exponent of $1.02$, consistent with the
$1/\gamma$ law; the full-sweep exponent is $0.71$. At finite amplitude the response departs from the linear
reference and saturates, while the driven endpoint remains attracting.
The mode tracking, amplitude checks, and their numerical ranges accompany
the coupling curve in the supplement (\S\ref{supp:ping-coupling-results}).
Thus JR and PING share the critical mode's $1/\gamma$ gain law despite
their different detectors; the law does not describe every mode or the
entire finite-amplitude sweep.

\ifjne\else
\subsection{Mode competition and finite-amplitude gain in PING}\label{supp:ping-coupling-results}
In NMM2 PING, $J$ is the mean inter-neuron QIF synaptic weight.
Sweeping $J$ at fixed background drive toward the gamma Hopf reveals two
competing modes (Fig.~\ref{fig:nmm2_jcurve}). A second mode near $55$~Hz carries
the linear-response peak over $14.6\lesssim J\lesssim16.2$; the critical mode
carries it at lower couplings and again near the Hopf, where its frequency is
about $50$~Hz. In that near-Hopf interval,
$0.0014\lesssim\gamma\lesssim0.0081$~ms$^{-1}$, the linear gain scales as
$1/\gamma$ with a fitted exponent of $1.02$ over a $5.8$-fold damping range.
The full-sweep exponent is $0.71$, so a single-mode law does not describe the
whole curve. Coupling-dependent sensitivity \cite{clusella2023} therefore
includes changes in which mode carries the peak.

At $A_f=8$, the simulated response exceeds the normalized linear reference by
up to $48\%$ in the amplitude check, then saturates. At a quarter of that
amplitude ($A_f=2$, one sixteenth of the power), the departures stay within
$11\%$ at the nine tested couplings. The driven endpoint remains attracting under both waveforms
(\S\ref{app:ping-stability}).
The alpha JR and gamma NMM2 circuits share the $1/\gamma$ gain scaling of their critical mode despite their different
mixers, with amplification adjustable through coupling.

\fi

\begin{figure}[t!]
\centering
\includegraphics[width=\textwidth]{fig_qif_raster_twotone.pdf}
\caption{\textbf{Beat-locked activity superposed on gamma.}
QIF network, $N=8000$ per population, $A_f=35$ per tone, carrier $300$~Hz.
Black ticks: 400 E neurons; blue: smoothed rate on its own scale; red:
envelope in arbitrary units. The last $200$~ms are shown.
\emph{(a,b)} Field off/on at $\bar\eta=0.85$, beat $55$~Hz: induced gamma
near $52$~Hz dominates. \emph{(c,d)} Field off/on at $\bar\eta=1.5$,
beat $42$~Hz: autonomous gamma remains dominant.
Right-hand panels fold rate on beat phase over $44$ and $34$ cycles, respectively,
normalized to each condition's mean (off gray, on blue; scale 0--2).
Folding suppresses non-beat-locked gamma and exposes the smaller beat component;
it does not establish gamma-amplitude gating or phase capture.}
\label{fig:qif_raster}
\end{figure}

\subsection{QIF responses distinguish beat modulation from field-induced gamma}\label{sec:qif}
The finite QIF network shows beat-locked modulation at the lower tested
amplitude and field-induced gamma at the higher amplitude. Its $8000$-neuron populations test
these regimes at the spike level
\cite{montbrio2015,ruffini2025rosetta,clusella2023}
(Figs.~\ref{fig:qif_raster},~\ref{fig:qif_regimes}).

At $\bar\eta=0.85$ and $A_f=10$, the original finite-cutoff run gives
beat depth $0.09$ and a $2.8\%$ mean-rate increase. Its quantitative depth
remains unconverged because of cutoff and population sensitivity (supplement,
\S\ref{app:qif-checks}--\ref{supp:theta-pilot}). The estimated rectified
E-drive shift remains below $\delta\eta_E^*$, the shift in excitatory drive
required to cross the Hopf while holding the inhibitory drive fixed.
At $\bar\eta=1.5$, the original finite-cutoff run gives a $42$~Hz depth
of about $0.02$ and a $1.6\%$ mean-rate increase, while ongoing gamma dominates.

At $A_f=35$ and $\bar\eta=0.85$, a line near $52$~Hz dominates
(depth $0.70$), alongside the $55$~Hz beat (depth $0.29$). This is a
finite-amplitude example: carrier averaging predicts that the rectified
E-drive shift exceeds $\delta\eta_E^*$ and destabilizes the averaged equilibrium.
Direct mean-field integration reproduces the off-beat line, identifying the
dominant response as field-induced gamma (Fig.~\ref{fig:qif_raster}b).
The drive and stability estimates at both amplitudes are given in the
supplement (\S\ref{app:qif-estimate}).

At $\bar\eta=1.5$, autonomous gamma near $53$--$54$~Hz remains about
four times larger than the $42$~Hz beat component (Fig.~\ref{fig:qif_raster}d).
Their coexistence produces slower burst-height variations. Mixing produces
sidebands near $95$ and $11$~Hz, as in endogenous amplitude modulation
\cite{ruffini2025ham}; these establish nonlinear mixing without distinguishing
gamma-amplitude gating from phase modulation (Fig.~\ref{fig:qif_timing}b).
Field-free gamma frequencies agree between the finite network and mean field
(Fig.~\ref{fig:qif_timing}a).

To separate steady induction from beat-dependent activity, we compare TI
with an unmodulated carrier matched in RMS input current, in model units
(Fig.~\ref{fig:qif_regimes}c,d). In the mean field, their rectified E-drive
shifts agree within $3\%$ and their mean-rate increases are similar.
At $A_f=35$, both induce gamma near $51$~Hz; TI also produces the
$55$~Hz beat. At $A_f=10$, neither induces gamma, and TI alone produces
a resolved beat response. This control has been tested only in the mean
field; its spectral-leakage check is in the supplement
(\S\ref{supp:qif-protocol}).

\begin{figure}[t!]
\centering
\includegraphics[width=\textwidth]{fig_qif_regimes.pdf}
\caption{\textbf{TI and an equal-RMS carrier across the gamma threshold.}
Background $\bar\eta=0.85$, carrier $300$~Hz, beat $55$~Hz.
\emph{(a,b)} Spiking network at $A_f=10$ and $35$ per tone.
Black: 400 of 8000 E neurons; blue: population rate on a common scale;
red: envelope. Panels below the rasters show mean-normalized phase folds
(off gray, on blue).
Beat depths at this cutoff are about $0.09$ and $0.3$, with mean-rate increases of $2.8\%$
and $10.5\%$; the larger input also induces gamma.
\emph{(c,d)} Mean-field spectra under TI (blue) and one carrier of peak
amplitude $\sqrt2 A_f$ (red), normalized to mean rate. Both induce gamma
at the larger input; TI adds the $55$~Hz beat (dashed).
Hann windows span $800$~ms after $800$~ms settling.
Spiking panels use $V_{\rm peak}=100$; weak-drive depth depends on cutoff and population size.
RMS matching is of input current in model units, not physical power.}
\label{fig:qif_regimes}
\end{figure}

At $A_f=35$, population beat modulation exceeds the relative mean-rate
change by about threefold. Across the two conditions, mean rate rises
$8$--$11\%$, while beat depth rises from below $0.003$ to $0.29$
($\bar\eta=0.85$) and $0.23$ ($\bar\eta=1.5$)
(Fig.~\ref{fig:qif_timing}c). Individual locking remains weak: median
pairwise phase consistency rises from $-0.02$ in the field-free control
to $0.003$--$0.005$, and pooled vector strength from below $0.02$ to
$0.12$--$0.13$. Many weakly biased cells therefore produce a population
beat component. These finite-amplitude results do not establish its size
under weak physiological stimulation.

